\chapter{Gen Z Đi Làm Sớm}
\markboth{Gen Z Đi Làm Sớm}{Gen Z Đi Làm Sớm}

\section{Muốn Tự Lập Từ Sớm}
\begin{truyen}{Muốn Tự Lập Từ Sớm}{Wanting Independence Early}
\chuhoa{H}{ọc sinh:} ``Em muốn đi làm sớm để tự lập, chứ em không muốn cứ xin tiền mãi.''
\textit{(Student: ``I want to start working early to become independent. I do not want to keep asking for money forever.'')}

\teacherpair{Khát vọng tự lập rất đáng quý. Nhưng tự lập không chỉ là có đồng tiền đầu tiên. Nó còn là biết giữ lời, chịu áp lực, làm việc với người khó, và không bỏ cuộc chỉ vì bị nhắc vài câu. Nhiều bạn mê cảm giác có lương hơn là mê quá trình trở thành người làm việc tử tế.}{``The desire for independence is admirable. But independence is not only about earning your first money. It also means keeping commitments, enduring pressure, working with difficult people, and not quitting after a few corrections. Many young people love the feeling of having income more than the process of becoming someone who works properly.''}
\end{truyen}

\section{Đi Làm Không Phải Để Đóng Kịch Người Lớn}
\begin{truyen}{Đi Làm Không Phải Để Đóng Kịch Người Lớn}{Working Is Not Performing Adulthood}
\chuhoa{H}{ọc sinh:} ``Đi làm sớm nghe ngầu hơn đi học suốt chứ thầy.''
\textit{(Student: ``Working early sounds cooler than just studying all the time, teacher.'')}

\teacherpair{Ngầu là thứ rẻ tiền lắm. Đi làm sớm không sai, nhưng nếu em vào đời chỉ để có ảnh check-in công sở, vài câu kể khổ và một cảm giác mình trưởng thành hơn người khác, thì em đang diễn vai người lớn chứ chưa sống được trách nhiệm người lớn.}{``Coolness is cheap. Working early is not wrong. But if you enter the workforce only for office check-ins, a few hardship stories, and the feeling of being more mature than others, then you are performing adulthood, not carrying its responsibility.''}
\end{truyen}

\section{Đi Sớm Hay Muộn, Cũng Phải Có Nền}
\begin{truyen}{Đi Sớm Hay Muộn, Cũng Phải Có Nền}{Early or Late, Work Still Needs Foundation}
\chuhoa{H}{ọc sinh:} ``Vậy khi nào đi làm sớm là tốt?''
\textit{(Student: ``Then when is working early actually good?'')}

\teacherpair{Khi em biết mình đi để học kỷ luật, học va chạm, học giá trị lao động, chứ không phải để trốn học hay chứng minh mình hơn người. Và dù đi sớm hay muộn, em vẫn cần một cái nền: kỹ năng, tác phong, khả năng học tiếp, và lòng tự trọng đủ để không bán rẻ mình chỉ vì muốn có tiền nhanh.}{``When you know you are going in order to learn discipline, friction, and the value of labor, not to escape study or prove superiority. And whether you work early or late, you still need a foundation: skill, work ethic, capacity to keep learning, and enough self-respect not to sell yourself cheaply for fast money.''}
\end{truyen}

\begin{baihoc}
Đi làm sớm không tự động làm một người trưởng thành hơn.
\textit{(Starting work early does not automatically make someone more mature.)}

Điều quyết định vẫn là em bước vào lao động với tâm thế nào và chuẩn bị đến đâu.
\textit{(What matters is the mindset and preparation you bring into work.)}
\end{baihoc}

\begin{ghinhoanh}
\textbf{Short English to remember:}

Early work is not instant maturity.

Income matters less than inner readiness.
\end{ghinhoanh}

\begin{tuhocnhanh}
1. Em muốn đi làm sớm vì học hỏi thật, hay vì muốn thoát một điều mình đang ghét? \textit{(Do you want to work early to learn, or simply to escape something you dislike?)}

2. Nếu đi làm từ bây giờ, điểm yếu nào của em sẽ lộ rõ nhất? \textit{(If you started working now, which weakness of yours would show up first?)}
\end{tuhocnhanh}

\ngancach